% GERADO por tools/generation/gerar_secoes_latex.py a partir de
% artifacts/paper/secoes/01-introducao.md. Nao editar a mao: a proxima geracao
% desfaz. Corrija o Markdown e regere.
Organizações realizam sua estratégia por projetos, e a exigência de dar visibilidade a eles há muito deixou de ser interna: a Lei de Acesso à Informação estabeleceu o acesso à informação pública como regra, e não como exceção \cite{L12527}, e a literatura de gerenciamento passou a tratar a transparência como atributo do próprio desenvolvimento do projeto \cite{betta2018transparency}. Atender a essa exigência mobiliza \textbf{pessoas} --- as partes interessadas dentro e fora da organização executora ---, \textbf{processos e procedimentos} --- as rotinas de coleta, tratamento e divulgação da informação sobre os projetos --- e \textbf{tecnologias} --- a infraestrutura que sustenta essas rotinas. O observatório de projetos é a estrutura sociotécnica que articula os três: uma entidade dedicada à observação contínua de um conjunto de projetos, que centraliza, analisa e divulga informação sobre eles \cite{vieira2021model}.

Por que estruturas assim existem é o que a Organizational Information Processing Theory explica \cite{galbraith1973designing,galbraith1974organization}: quanto maior a incerteza de uma tarefa, maior a informação que precisa ser processada entre quem a executa e quem decide sobre ela, e a organização responde ampliando sua capacidade de processamento. Um observatório de projetos é um mecanismo desse tipo, com uma particularidade: a capacidade que ele amplia não é a de quem gerencia os projetos, e sim a de quem precisa acompanhá-los de fora --- financiadores, órgãos de controle e sociedade. É essa assimetria que o sustenta como categoria própria (\S{}2.1).

O domínio de aplicação deste artigo são os observatórios de projetos em organizações que prestam contas a terceiros --- universidades públicas, agências de fomento e órgãos executores de políticas ---, e seu modelo de referência é o Modelo para Observatórios de Projetos (MPO) \cite{vieira2021model,vieira2022thesis}, que organiza 61 conceitos ``distribuídos em três níveis hierárquicos'' e em três dimensões. O problema organizacional está em como esse modelo é descrito. Sendo prosa, o MPO não fixa a interpretação de seus conceitos: uma iniciativa que se declara aderente a ele não tem como demonstrar em que a aderência consiste, duas que usam o mesmo termo podem estar nomeando entidades distintas, e a prestação de contas --- que exige percorrer o caminho do conteúdo divulgado até a fonte que o originou --- não encontra no modelo o encadeamento que permitiria auditá-la.

A tese defendida aqui é que esses efeitos não são acidente de um modelo particular: \textbf{modelos conceituais de referência descritos em linguagem natural carregam deficiências representacionais sistemáticas e invisíveis a quem os aplica}. Invisíveis porque o texto não é contraditório: é subdeterminado, e a lacuna só aparece quando alguém precisa decidir o que cada conceito é. Submeter o MPO a uma análise ontológica fundamentada na Unified Foundational Ontology (UFO) tornou nove dessas lacunas observáveis, cada uma classificada pela tipologia da Representation Theory \cite{wand1993ontological}, corrigida com justificativa ontológica e rastreada até o ponto do modelo que a permitiu.

A contribuição declarada, portanto, não é formalizar o MPO --- uma formalização já existe, publicada por terceiros, e é ela o objeto analisado. A contribuição é a \textbf{análise ontológica} e o \textbf{modelo revisado} que dela resultou, com um procedimento reusável para auditar modelos de referência descritos em prosa e a evidência antes/depois que separa o que a análise afirma do que as ferramentas atestam. Esse resultado responde ao Desafio 3 do II GranDSI-Br, \textit{Eco(Sistemas\textsuperscript{2}) de Informação} \cite{araujo2025grandsi}: observatórios são ecossistemas de informação que só se integram se os conceitos que trocam significarem o mesmo dos dois lados, e é essa interoperabilidade semântica que uma especificação verificável do modelo de referência torna possível. Secundariamente, alinha-se ao Desafio 2, \textit{SI Inteligentes sob a Perspectiva Sociotécnica}, ao explicitar a camada de agentes que a formalização tratava como tipo único.

O artigo segue assim. A \S{}2 apresenta os observatórios de projetos, o MPO, a Representation Theory e a UFO; a \S{}3 situa o trabalho frente à literatura e a \S{}4 descreve o desenho de pesquisa. A \S{}5 apresenta a análise ontológica e as nove deficiências, com seus traços e correções; a \S{}6, a OntoMPO revisada, e a \S{}7, sua avaliação. A \S{}8 discute lições e limitações e a \S{}9 conclui.
